\documentclass{article}
\usepackage{PRIMEarxiv} % Core package for the PrimeAI Template
\usepackage{amsmath, amssymb, amsthm, bm, dcolumn} % Add amsthm here for the proof environment
\usepackage[numbers,sort&compress]{natbib} % Natbib for citations
\usepackage{graphicx} % For high-quality images
\usepackage{multicol} % Multi-column figures/tables
\usepackage{hyperref} % Hyperlinks
\usepackage{listings} % Code listings
\usepackage{authblk} % For structured affiliations
% Define theorem style (optional)
\theoremstyle{plain}
\newtheorem{theorem}{Theorem}
\renewcommand{\qedsymbol}{}

% Custom commands for primes and qubit encoding
\newcommand{\primeQubit}[1]{|p_{#1}\rangle}
\newcommand{\primeGate}[1]{U_{p_{#1}}}

\usepackage{wrapfig}
\usepackage[pscoord]{eso-pic}
\usepackage[fulladjust]{marginnote}
\reversemarginpar

% Typesetting improvements without footnote patching
\usepackage[protrusion=true, expansion=true, tracking=false]{microtype}
\microtypecontext{spacing=nonfrench}

% Line numbers
\usepackage[right]{lineno}

% Text layout - adjust as needed
\raggedright
\setlength{\parindent}{0.5cm}
\textwidth 5.25in 
\textheight 8.75in

% Set double spacing
\usepackage{setspace} 
\doublespacing

% Adjust width for specific content
\usepackage{changepage}

% Adjust caption style
\usepackage[aboveskip=1pt,labelfont=bf,labelsep=period,singlelinecheck=off]{caption}

% Remove brackets from references
\makeatletter
\renewcommand{\@biblabel}[1]{\quad#1.}
\makeatother

% Header, footer, and page numbers
\usepackage{lastpage,fancyhdr}
\pagestyle{fancy}
\fancyhf{}
\fancyhead[L]{Preprint - PrimeAI Enhanced Template}
\fancyfoot[C]{\scriptsize Multiplicity Theory © 2025 Citizen Gardens \\ Licensed Under MIT and CC BY-NC-SA 4.0.}
\fancyfoot[R]{Page \thepage\ of \pageref{LastPage}}
\renewcommand{\footrule}{\hrule height 2pt \vspace{2mm}}

\begin{document}

\title{Classical \& Quantum Hall Effect \\ with Multiplicity Theory}

\author{Ryan O. Van Gelder}
\author{Nicholas Galioto}
\affil{Citizen Gardens - The Foundation of Multiplicity\\info@citizengardens.org}
\date{\today}

\maketitle
\section{Classical Hall Effect in a Multiplicative Framework}

The classical Hall Effect describes the emergence of a transverse voltage when a conductor carrying current is exposed to a perpendicular magnetic field. Mathematically, it is governed by:

\begin{equation}
\mathbf{F} = q (\mathbf{E} + \mathbf{v} \times \mathbf{B})
\end{equation}

where:
\begin{itemize}
\item $\mathbf{F}$ is the Lorentz force,
\item $q$ is the charge of the carrier,
\item $\mathbf{E}$ is the electric field,
\item $\mathbf{v}$ is the drift velocity of charge carriers,
\item $\mathbf{B}$ is the applied magnetic field.
\end{itemize}

\subsection{Multiplicative Tensor Representation of Hall Voltage}
Instead of treating the Hall voltage $V_H$ as an additive function, we introduce a multiplicative transformation:

\begin{equation}
V_H = \frac{B}{n q d} I
\end{equation}

where:
\begin{itemize}
\item $n$ is the charge carrier density,
\item $d$ is the thickness of the material,
\item $I$ is the current.
\end{itemize}

Using tensor notation, the Hall voltage can be rewritten as:

\begin{equation}
V_H = \mathbf{T_H} \cdot \mathbf{J}
\end{equation}

where $\mathbf{T_H}$ is a multiplicative transformation tensor encoding the interaction between charge flow, material properties, and the external magnetic field.

\section{Multiplicative Modeling of Carrier Dynamics}

In classical approaches, charge carrier dynamics are modeled additively. However, Multiplicity Theory suggests that charge transport in Hall systems should be modeled using multiplicative differential equations.

\subsection{Multiplicative Drift Velocity Equation}
Using standard drift velocity:

\begin{equation}
\mathbf{v_d} = \frac{\mathbf{E} \times \mathbf{B}}{B^2}
\end{equation}

we introduce a multiplicative form:

\begin{equation}
\mathbf{v_d} = \mathcal{M}(\mathbf{E}, \mathbf{B}) \cdot \mathbf{E}
\end{equation}

where $\mathcal{M}(\mathbf{E}, \mathbf{B})$ is a multiplicative operator encoding field interactions.

\section{Quantum Hall Effect in a Multiplicative Framework}

For the Quantum Hall Effect (QHE), conductivity is quantized:

\begin{equation}
\sigma_H = \frac{n e^2}{h}
\end{equation}

We reformulate this using multiplicative eigenfunctions:

\begin{equation}
\sigma_H = \lambda_H \cdot f(n, e, h)
\end{equation}

where $\lambda_H$ is a multiplicative eigenvalue representing topological quantization.

\subsection{Multiplicative Topology Transformations}
QHE plateaus can be described via a multiplicative topology operator:

\begin{equation}
\mathbf{T_Q} \cdot \sigma_H = k \cdot \frac{e^2}{h}
\end{equation}

where $\mathbf{T_Q}$ acts as a multiplicative topological transformation tensor.

\section{Computational Simulations Using Multiplicative Networks}

To simulate Hall conductivity under multiplicative models, we define a multiplicative neural network where:

\begin{equation}
\sigma_{H}^{(i+1)} = W_H^{(i)} \cdot \sigma_H^{(i)} + B_H^{(i)}
\end{equation}

where:
\begin{itemize}
\item $W_H^{(i)}$ is a multiplicative weight matrix,
\item $B_H^{(i)}$ is an additive correction term for numerical stability.
\end{itemize}

This allows adaptive modeling of Hall resistivity in different materials.

\section{Experimental Design Based on Multiplicative Analysis}
To test these models, we propose an experiment where:
\begin{enumerate}
\item Variable Magnetic Fields – Measure Hall voltage at different field strengths to observe multiplicative scaling laws.
\item Temperature Modulation – Analyze how thermal effects multiplicatively alter charge mobility.
\item Dynamically Tuned Conductors – Use materials where carrier density is externally tunable to validate multiplicative models.
\end{enumerate}

\section{Prime-Encoding the Hall Effect}

Prime encoding provides an alternative computational method to analyze the Hall Effect using prime numbers. The fundamental concept is to map electrical and magnetic properties to a prime-indexed function space, enabling a discrete analysis of continuous interactions.

\subsection{Prime Representation of Charge Carrier Density}
Instead of using a continuous function for charge carrier density $n$, we represent it as a prime sequence:

\begin{equation}
n = p_k, \quad \text{where } p_k \text{ is the } k\text{th prime number.}
\end{equation}

This allows for a discrete, non-uniform sampling approach, enhancing numerical stability in quantum simulations.

\subsection{Prime-Modulated Conductivity}
Conductivity $\sigma_H$ can be restructured as:

\begin{equation}
\sigma_H = \frac{p_m e^2}{h}
\end{equation}

where $p_m$ is a prime-selected modulation factor, adjusting for material properties dynamically.

\section{Experimental Design Based on Multiplicative Analysis}
To test these models, we propose an experiment where:
\begin{enumerate}
\item Variable Magnetic Fields – Measure Hall voltage at different field strengths to observe multiplicative scaling laws.
\item Temperature Modulation – Analyze how thermal effects multiplicatively alter charge mobility.
\item Dynamically Tuned Conductors – Use materials where carrier density is externally tunable to validate multiplicative models.
\end{enumerate}

\section{Conclusion}
Using Multiplicity Theory, we introduce a new way to study the Hall Effect, emphasizing multiplicative tensor modeling, computational networks, and quantum transformations. Additionally, we incorporate prime-encoded formulations for discrete representations, which can enhance numerical stability and computational efficiency.
\end{document}